\chapter{Технологическая часть}

\section{Средства разработки}
В качестве языка программирования со статической типизацией был выбран язык C++ (стандарт C++20). Для массива использовался контейнер стандартной библиотеки \texttt{std::vector<int>}.

\section{Реализация алгоритмов}

На рисунке ~\ref{fig:brute_force_search} представлена программная реализация алгоритма поиска полным перебором.
\begin{figure}[H]
	\centering
		\lstinputlisting[
		firstline=6, lastline=19
		]{algorithms.hpp}
		\caption{Поиск полным перебором}
		\label{fig:brute_force_search}
\end{figure}

На рисунках~\ref{fig:bin_search_1} и~\ref{fig:bin_search_2} представлены два варианта программных реализаций нерекурсивного алгоритма бинарного поиска.
\begin{figure}[H]
	\centering
		\lstinputlisting[
		firstline=21, lastline=41
		]{algorithms.hpp}
		\caption{Нерекурсивный бинарный поиск с отложенным выходом}
		\label{fig:bin_search_1}
\end{figure}

\begin{figure}[H]
	\centering
		\lstinputlisting[
		firstline=43, lastline=60
		]{algorithms.hpp}
		\caption{Нерекурсивный бинарный поиск с предварительным выходом}
		\label{fig:bin_search_2}
\end{figure}

На рисунках ~\ref{fig:bin_search_r} и~\ref{fig:wrapper} представлены программная реализация рекурсивного алгоритма бинарного поиска и соответствующая функция-обёртка.
\begin{figure}[H]
	\centering
		\lstinputlisting[
			firstline=62, lastline=75
		]{algorithms.hpp}
		\caption{Рекурсивный бинарный поиск}
		\label{fig:bin_search_r}
\end{figure}

\begin{figure}[H]
	\centering
		\lstinputlisting[
			firstline=77, lastline=82
		]{algorithms.hpp}
		\caption{Оберточная функция для рекурсивного бинарного поиска}
		\label{fig:wrapper}
\end{figure}

\section{Тестирование реализаций}
В таблице~\ref{tbl:tests} представлены результаты функционального тестирования разработанных реализаций алгоритмов поиска.

\begin{table}[H]
	\caption{Функциональные тесты}
	\label{tbl:tests}
	\centering
	\begin{tblr}{
			width = \linewidth,
			colspec = {Q[c,m] Q[c,m] Q[c,m] Q[c,m] Q[c,m] X[1.6,c,m]},
			hlines, vlines,
			row{1,2} = {font=\bfseries},
		}
		\SetCell[r=2]{c} №
		& \SetCell[c=2]{c} Входные данные &
		& \SetCell[c=2]{c} Выходные данные &
		& \SetCell[r=2]{c} Комментарий \\
		& Массив & $x$ & Ожидаемые & Полученные & \\
		1  & [10]                                    & 10   & 0         & 0         & Один элемент, совпадение \\
		2  & [10]                                    & 5    & $-1$      & $-1$      & Один элемент, не найден \\
		3  & [10, 20]                                & 10   & 0         & 0         & Два элемента, первый \\
		4  & [10, 20]                                & 20   & 1         & 1         & Два элемента, второй \\
		5  & [10, 20]                                & 15   & $-1$      & $-1$      & Два элемента, не найден \\
		6  & [1, 3, 5, 7, 9]                         & 1    & 0         & 0         & Первый элемент, нечётная длина \\
		7  & [1, 3, 5, 7, 9]                         & 9    & 4         & 4         & Последний элемент, нечётная длина \\
		8  & [1, 3, 5, 7, 9]                         & 5    & 2         & 2         & Середина, нечётная длина \\
		9  & [1, 3, 5, 7, 9]                         & 0    & $-1$      & $-1$      & Не найден, нечётная длина \\
		10 & [1, 3, 5, 7]                            & 1    & 0         & 0         & Первый элемент, чётная длина \\
		11 & [1, 3, 5, 7]                            & 7    & 3         & 3         & Последний элемент, чётная длина \\
		12 & [1, 3, 5, 7]                            & 5    & 2         & 2         & Середина, чётная длина \\
		13 & [1, 3, 5, 7]                            & 0    & $-1$      & $-1$      & Не найден, чётная длина \\
	\end{tblr}
\end{table}

Все разработанные варианты реализации алгоритма поиска успешно прошли функциональное тестирование.

\section{Вывод}
В результате технологической части были выбраны язык программирования и средства разработки, представлены листинги исходного кода разработанных алгоритмов поиска, выполнено функциональное тестирование, результаты которого подтвердили корректность работы всех реализаций.